% Propietario conservado: /Users/ruben/Documents/New project/output/EXTREMA_MEDIA_RAZON_RECIPROCIDAD_ES_EN_20260918/ARTICULO/ES/source/sections/iv_accion_antecedente.tex
% RESULTADO_RECUPERADO. II REV10, 05_accion.tex:1--195,
% revision_planck.tex:13--93 y definicion_planck_rev06.tex:4--34.
% Selección: antecedentes de acción de la elipse; no incorpora aplicaciones
% electrónicas ni el contraste metrológico, que no determinan el radio.
\section{Estructura determinantal y secciones orientadas de acción}
\label{x_reciprocidad:iv:action:section}

Las publicaciones de \(\pi,\varphi,e,\alpha\) proceden del estado
APP--TRIT--TPK de las secciones
\ref{x_reciprocidad:sec:nucleo}--\ref{x_reciprocidad:sec:alpha}. El registro dodecafásico de
\ref{x_reciprocidad:sec:k-registro} conserva su orientación, sus residuos y
cocientes y la memoria del retorno
\eqref{x_reciprocidad:eq:holonomia-memoria}. Sobre ese registro se definen dos operaciones
sucesivas: una norma local determina la valoración de la sección de
acción, y un carácter con continuación regional determina sus
coordenadas anterior y posterior al retorno. Ambas operaciones
transmiten a la construcción angular sus reglas y normalizaciones.

\subsection{Cociente local y norma determinantal}
\label{x_reciprocidad:iv:action:determinantal}

El bloque local de paso tres del registro dodecafásico, identificado
en \ref{x_reciprocidad:subsec:k-hadamard}, es
\begin{equation}
 H_4=\begin{pmatrix}
 1&1&1&1\\
 1&1&-1&-1\\
 1&-1&1&-1\\
 1&-1&-1&1
 \end{pmatrix},
 \qquad G_H=\mathbb Z^4/H_4\mathbb Z^4.
 \label{x_reciprocidad:iv:action:h4}
\end{equation}
El cociente corresponde a este bloque local. Los tres bloques del
registro completo conservan sus posiciones; la operación de norma
que sigue actúa sobre las hojas de \(G_H\).

\begin{lemma}[Estructura del cociente local]
\label{x_reciprocidad:iv:action:smith}
La forma normal de Smith de \(H_4\) es
\(\operatorname{diag}(1,2,2,4)\). Por consiguiente,
\[
 G_H\simeq\mathbb Z/2\mathbb Z\oplus\mathbb Z/2\mathbb Z
       \oplus\mathbb Z/4\mathbb Z,\qquad |G_H|=16.
\]
\end{lemma}
\begin{proof}
El máximo común divisor de las entradas es uno. Los menores de
orden dos son pares y uno tiene valor \(-2\), por lo que su máximo
común divisor es dos. Las identidades
\(H_4H_4^{\mathsf T}=4I\) y \(\det H_4=-16\) muestran que los
cofactores valen \(\pm4\); el tercer divisor determinantal es
cuatro. El cuarto es \(16\). Los cocientes sucesivos de
\(1,2,4,16\) son \(1,2,2,4\); su producto determina el orden
del cociente.
\end{proof}

\begin{definition}[Lector determinantal de acción]
\label{x_reciprocidad:iv:action:norma}
La sección escalar \(\alpha\), constante sobre las hojas de \(G_H\),
y una aplicación de la autoescala \(\varphi\) determinan
\begin{equation}
 \Sigma_H=\varphi\,\mathrm N_{G_H}(\alpha),\qquad
 \mathrm N_{G_H}(\alpha)=\prod_{g\in G_H}\alpha=\alpha^{16}.
 \label{x_reciprocidad:iv:action:sigma}
\end{equation}
Su valoración decimal es
\[
 \operatorname{ord}_{10}(\Sigma_H)
   =\lfloor\log_{10}\Sigma_H\rfloor.
\]
\end{definition}
La potencia dieciséis procede del número de hojas una vez fijada
la norma multiplicativa. La norma y la aplicación única de
\(\varphi\) forman parte del lector declarado: el orden del grupo,
considerado aisladamente, admite otros funcionales.

\begin{theorem}[Valoración decimal de la sección de acción]
\label{x_reciprocidad:iv:action:decada}
El lector anterior, aplicado a las coordenadas HMT ya generadas,
satisface
\begin{equation}
 10^{-34}<\Sigma_H<10^{-33},
 \qquad \operatorname{ord}_{10}(\Sigma_H)=-34.
 \label{x_reciprocidad:iv:action:orden-decimal}
\end{equation}
\end{theorem}
\begin{proof}
Se utilizan las inclusiones racionales de las publicaciones
anteriores,
\begin{equation}
 \frac{1618}{1000}<\varphi<\frac{1619}{1000},
 \qquad
 \frac{7297}{10^6}<\alpha<\frac{7298}{10^6}.
 \label{x_reciprocidad:iv:action:intervalos}
\end{equation}
La función \((x,a)\mapsto xa^{16}\) es creciente en ambas
variables positivas. Las desigualdades enteras
\[
 1618\cdot7297^{16}>10^{65},
 \qquad 1619\cdot7298^{16}<10^{66}
\]
dan
\[
 10^{-34}
 <\frac{1618}{1000}\left(\frac{7297}{10^6}\right)^{16}
 <\Sigma_H
 <\frac{1619}{1000}\left(\frac{7298}{10^6}\right)^{16}
 <10^{-33}.
\]
La definición de la valoración concluye la prueba. Las inclusiones
son controles posteriores de las coordenadas ya producidas; la
operación determinantal precede a cualquier carta de unidades.
\end{proof}

\subsection{Carácter de acción y continuación regional}
\label{x_reciprocidad:iv:action:caracter}

La carta central conserva la matriz
\[
 B_c=\begin{pmatrix}7&2\\2&7\end{pmatrix},
 \qquad\lambda_+=9,\quad\lambda_-=5.
\]
El calendario tiene longitud \(12\), la órbita del paso tres tiene
longitud cuatro y el cociente censal utilizado es
\(104976/1944=54\), con los censos definidos en la
sección~\ref{x_reciprocidad:sec:extension}. Estos invariantes determinan los coeficientes
asignados a los grados \(2,3,4,5\) del carácter:
\begin{equation}
 \frac9{16}=\frac{\lambda_+}{4^2},\qquad
 \frac59=\frac{\lambda_-}{\lambda_+},\qquad
 \frac7{48}=\frac{\operatorname{tr}(B_c)/2}{4\cdot12},
 \qquad
 \frac1{54}=\frac{1944}{104976}.
 \label{x_reciprocidad:iv:action:coeficientes}
\end{equation}
La asignación a esos grados y la alternancia de signos pertenecen
a la definición del carácter analítico. El contenido operativo
reúne el soporte, la orientación y esa regla de asignación.

\begin{definition}[Carácter local y lector regional de retorno]
\label{x_reciprocidad:iv:action:lectores}
Sobre las publicaciones \(\alpha,\varphi,\pi\), el carácter de
quinto orden, la transferencia regional y su continuación son
\begin{align}
 H_5={}&\frac12\sqrt{\frac{1000\alpha}{\varphi}}-\alpha
       +\frac9{16}\alpha^2-\frac59\alpha^3
       +\frac7{48}\alpha^4-\frac1{54}\alpha^5,
       \label{x_reciprocidad:iv:action:h5}\\
 D_A={}&\exp\!\left(-\frac{100\pi\alpha}{9}\right),
       \label{x_reciprocidad:iv:action:da}\\
 \mathcal C_\pi(\alpha)={}&169\alpha^6+
       \frac{D_A\alpha^7}{1-D_A\alpha},
       \label{x_reciprocidad:iv:action:cpi}\\
 \eta_{\mathrm{ret}}={}&H_5-\frac{90}{\pi}\mathcal C_\pi(\alpha).
       \label{x_reciprocidad:iv:action:eta-ret}
\end{align}
El entero \(169\) es el selector regional de las publicaciones
construidas en \ref{x_reciprocidad:sec:generacion}: se conserva la emisión regional
\(169\,|\,272\,|\,272\), distinguida de
\(418\,|\,674\,|\,521\) en la construcción excepcional.
La primera prolongación estricta de su continuación
tiene peso \(D_A\) y cada concatenación multiplica el peso por
\(D_A\). La normalización inicial es unitaria.
\end{definition}

La continuación se formula antes de evaluar en \(z=\alpha\).
Con \(D_A\) ya producido por su lector, escribimos
\(\mathscr C(z)=169z^6+\mathscr R(z)\), donde
\(\mathscr R(z)\in z^7\mathbb R[[z]]\).
La concatenación aumenta el grado en una unidad; el peso inicial
y la regla de transferencia dan
\begin{equation}
 \mathscr R(z)=D_Az^7+D_Az\mathscr R(z).
 \label{x_reciprocidad:iv:action:renovacion}
\end{equation}

\begin{proposition}[Unicidad de la continuación normalizada]
\label{x_reciprocidad:iv:action:unicidad-renovacion}
La ecuación \eqref{x_reciprocidad:iv:action:renovacion} determina una única serie:
\begin{equation}
 \mathscr C(z)=169z^6+\frac{D_Az^7}{1-D_Az}
             =169z^6+\sum_{n=7}^{\infty}D_A^{n-6}z^n.
 \label{x_reciprocidad:iv:action:serie-renovacion}
\end{equation}
Su realización es analítica para \(|D_Az|<1\).
\end{proposition}
\begin{proof}
El término independiente de \(1-D_Az\) es uno, de modo que ese
elemento es invertible en el anillo de series formales.
Resolver \eqref{x_reciprocidad:iv:action:renovacion} da la expresión racional.
Equivalentemente, el coeficiente de grado siete es \(D_A\) y cada
coeficiente posterior es \(D_A\) veces el anterior. La serie
geométrica converge absolutamente en el disco indicado.
\end{proof}

La normalización del primer peso interviene en la unicidad.
El impulso de grado seis y la razón de concatenación, sin ese
peso inicial, también son compatibles con
\[
 169z^6+\lambda\frac{D_Az^7}{1-D_Az}.
\]
La regla unitaria fija \(\lambda=1\) antes de la evaluación.
Para \(0<\alpha<1\), se tiene \(0<D_A<1\) y \(D_A\alpha<1\);
por tanto,
\begin{equation}
 \frac{D_A\alpha^7}{1-D_A\alpha}
   =\sum_{n=0}^{\infty}D_A^{n+1}\alpha^{n+7}.
 \label{x_reciprocidad:iv:action:cola}
\end{equation}
Después del término de índice \(N\), el resto exacto es
\(D_A^{N+2}\alpha^{N+8}/(1-D_A\alpha)\). La expresión racional
conserva así la continuación completa a cualquier precisión.

\subsection{Secciones positivas y acción por vuelta}
\label{x_reciprocidad:iv:action:secciones}

Sea \(\mathcal L_S\) la recta orientada de acción y
\(\mathcal U_S\) una base positiva. Sus secciones anterior y
posterior al retorno son
\begin{equation}
 \hbar_{\mathrm{pre}}=H_5\,10^{-34}\mathcal U_S,\qquad
 \hbar_{\mathrm{ret}}=\eta_{\mathrm{ret}}\,10^{-34}\mathcal U_S.
 \label{x_reciprocidad:iv:action:hbar}
\end{equation}
El retorno conserva la compensación regional
\(90\mathcal C_\pi(\alpha)/\pi\) y la razón
\begin{equation}
 \delta_{\mathrm{ret}}=H_5-\eta_{\mathrm{ret}},\qquad
 R_{\mathrm{act}}=\frac{\hbar_{\mathrm{pre}}}
                        {\hbar_{\mathrm{ret}}}
                =\frac{H_5}{\eta_{\mathrm{ret}}}.
 \label{x_reciprocidad:iv:action:razon}
\end{equation}

\begin{proposition}[Positividad y covariancia de la razón de acción]
\label{x_reciprocidad:iv:action:positividad}
En los intervalos \eqref{x_reciprocidad:iv:action:intervalos},
\(0<\eta_{\mathrm{ret}}<H_5\), y \(R_{\mathrm{act}}>1\).
La razón permanece fija al cambiar la base positiva común
\(\mathcal U_S\).
\end{proposition}
\begin{proof}
Los términos de \(\mathcal C_\pi\) son positivos. Para acotar
\(H_5\), bastan \(\alpha>0.007\), \(\alpha<0.008\) y
\(\varphi<1.619\). El término de raíz de \eqref{x_reciprocidad:iv:action:h5}
es mayor que uno y la suma de sus términos negativos es menor
que \(0.008001\); luego \(H_5>0.99\).
Con \(D_A<1\) y \(\pi>3\),
\[
 0<\frac{90}{\pi}\mathcal C_\pi
 <30\left(169(0.008)^6+
             \frac{(0.008)^7}{1-0.008}\right)<10^{-6}.
\]
Así, \(\eta_{\mathrm{ret}}>0.989999>0\) y
\(\eta_{\mathrm{ret}}<H_5\). El cambio de base divide ambas
coordenadas de acción por el mismo factor positivo; éste se
cancela en su cociente.
\end{proof}

La determinación estructural de la acción angular es el par
\((\hbar_{\mathrm{pre}},\hbar_{\mathrm{ret}})\) con su
identificación por retorno. Conserva la valoración de la norma
local, el carácter de acción y la continuación regional.
En cada sección \(s\in\{\mathrm{pre},\mathrm{ret}\}\), la acción
de una vuelta completa es
\begin{equation}
 h_s=2\pi\hbar_s.
 \label{x_reciprocidad:iv:action:vuelta}
\end{equation}
La carta posterior \(\mathcal U_S\mapsto\mathrm{J\,s}\) expresa
esas secciones en unidades de acción. El radio normalizado que
se construirá conserva su independencia respecto de esta elección
común de unidades.
